\section*{अध्याय \ref{ch:g5:raw-materials-and-recycling} --- कच्चे माल से पुनर्चक्रण तक}

\begin{solution}{exo:g5:raw-materials-and-recycling:1}
प्राकृतिक: लकड़ी, ऊन, पत्थर। निर्मित: काँच, स्टील, प्लास्टिक।
\end{solution}

\begin{solution}{exo:g5:raw-materials-and-recycling:2}
कागज़: लकड़ी (पेड़)। काँच: रेत (सोडा और चूना पत्थर के साथ)। सूत:
कपास का पौधा। प्लास्टिक: कच्चा तेल।
\end{solution}

\begin{solution}{exo:g5:raw-materials-and-recycling:3}
वह चट्टान जिससे कोई धातु निकाली जाती है। ऐलुमिनियम बॉक्साइट से आता है।
\end{solution}

\begin{solution}{exo:g5:raw-materials-and-recycling:4}
मर्तबान काँच के कूड़ेदान में, अख़बार कागज़ और गत्ते के साथ, डिब्बा
धातु के साथ, बोतल प्लास्टिक की बोतलों के साथ।
\end{solution}

\begin{solution}{exo:g5:raw-materials-and-recycling:5}
नए कच्चे माल की: प्रकृति से ली गई रेत, \omterm{def:g5:raw-materials-and-recycling:raw-material}{अयस्क} या तेल की जगह पुनर्चक्रित
सामग्री का उपयोग होता है।
\end{solution}

\begin{solution}{exo:g5:raw-materials-and-recycling:6}
नवीकरणीय: लकड़ी (अगर फिर से पेड़ लगाए जाएँ), ऊन, कपास। नवीकरणीय नहीं: कच्चा तेल,
लौह \omterm{def:g5:raw-materials-and-recycling:raw-material}{अयस्क}।
\end{solution}

\begin{solution}{exo:g5:raw-materials-and-recycling:7}
रेत से काँच बनाना एक \omterm{def:g5:irreversible-changes:chemical-change}{रासायनिक परिवर्तन} है: कुछ नया बनता है, और वापसी
का कोई सीधा रास्ता नहीं है।
\end{solution}

\begin{solution}{exo:g5:raw-materials-and-recycling:8}
उदाहरण के लिए: कम करना --- प्लास्टिक की पैकिंग के बिना फल ख़रीदना; फिर से उपयोग ---
काँच का मर्तबान फिर से भरना; \omterm{def:g5:raw-materials-and-recycling:recycling}{पुनर्चक्रण} --- पेय का डिब्बा धातु वाले कूड़ेदान में डालना।
\end{solution}

\begin{solution}{exo:g5:raw-materials-and-recycling:9}
हर सामग्री का \omterm{def:g5:raw-materials-and-recycling:recycling}{पुनर्चक्रण} अपने ढंग से होता है: किसी दूसरी सामग्री के टुकड़े
पिघले हुए काँच को ख़राब कर देंगे।
\end{solution}

\begin{solution}{exo:g5:raw-materials-and-recycling:10}
20 इकाइयों का बीसवाँ भाग: $20 \div 20 = 1$, यानी लगभग ऊर्जा की 1 इकाई।
\end{solution}

\begin{solution}{exo:g5:raw-materials-and-recycling:11}
\omterm{def:g5:raw-materials-and-recycling:recycling}{पुनर्चक्रण} में भी ऊर्जा, मशीनें और ढुलाई लगती है; जो वस्तु कभी बनती ही नहीं,
वह अपना \omterm{def:g5:raw-materials-and-recycling:raw-material}{कच्चा माल} और वह सारी ऊर्जा बचा लेती है।
\end{solution}
